% =====================================================================
%  Thème 5 — Angle et identités remarquables — FACILE — Exercices
% =====================================================================
\documentclass[11pt,a4paper]{article}
\usepackage[utf8]{inputenc}
\usepackage[T1]{fontenc}
\usepackage{lmodern}
\usepackage[french]{babel}
\usepackage{amsmath,amssymb,mathtools}
\usepackage{xcolor}
\usepackage{enumitem}
\usepackage{fancyhdr}
\usepackage[a4paper,margin=2.1cm,headheight=26pt,headsep=10pt]{geometry}

\definecolor{primary}{RGB}{222,70,80}
\definecolor{primarydark}{RGB}{150,30,42}

\pagestyle{fancy}\fancyhf{}
\fancyhead[L]{\small\textcolor{primarydark}{\textbf{Fiche 2} — Calcul vectoriel}}
\fancyhead[R]{\small\textcolor{primarydark}{\textsc{Thème 5 — Facile — Exercices}}}
\fancyfoot[C]{\small\thepage}
\renewcommand{\headrulewidth}{0.8pt}

\newcommand{\vc}[1]{\overrightarrow{#1}}
\providecommand{\norm}[1]{\left\lVert #1\right\rVert}
\newcommand{\exo}[1]{\medskip\noindent{\color{primary}\bfseries\large Exercice #1.}\par\nopagebreak\smallskip}

\begin{document}
\begin{center}
{\LARGE\bfseries\color{primarydark}Angle et identités remarquables — Niveau Facile}\\[2pt]
{\color{primary}\rule{0.55\linewidth}{1.1pt}}
\end{center}
\vspace{1mm}
{\small\itshape Objectif : calculer un cosinus, reconnaître un angle remarquable et appliquer
directement les identités remarquables.}
\vspace{2mm}

\exo{1} \textbf{Cosinus de l'angle.}
Calculer $\cos(\vec u,\vec v)$ (forme exacte simplifiée).
\begin{enumerate}[label=\textbf{\alph*)},leftmargin=2em,itemsep=2pt]
  \item $\vec u\,(1;0)$ et $\vec v\,(1;1)$.
  \item $\vec u\,(1;\sqrt3)$ et $\vec v\,(1;0)$.
  \item $\vec u\,(3;0)$ et $\vec v\,(0;5)$.
\end{enumerate}

\exo{2} \textbf{Angle remarquable.}
Donner la mesure (en degrés) de l'angle $\theta$ correspondant à chaque cosinus.
\begin{enumerate}[label=\textbf{\alph*)},leftmargin=2em,itemsep=2pt]
  \item $\cos\theta=\dfrac{\sqrt2}{2}$.
  \item $\cos\theta=\dfrac12$.
  \item $\cos\theta=\dfrac{\sqrt3}{2}$.
  \item $\cos\theta=0$.
\end{enumerate}

\exo{3} \textbf{Norme d'une somme (méthode directe).}
Calculer $\norm{\vec u+\vec v}$ en additionnant d'abord les composantes.
\begin{enumerate}[label=\textbf{\alph*)},leftmargin=2em,itemsep=2pt]
  \item $\vec u\,(1;2)$ et $\vec v\,(3;4)$.
  \item $\vec u\,(2;-1)$ et $\vec v\,(-1;3)$.
\end{enumerate}

\exo{4} \textbf{Identité remarquable.}
On donne $\vec u\,(1;1)$ et $\vec v\,(2;3)$.
\begin{enumerate}[label=\textbf{\alph*)},leftmargin=2em,itemsep=2pt]
  \item Calculer $\norm{\vec u}^2$, $\norm{\vec v}^2$ et $\vec u\cdot\vec v$.
  \item En déduire $\norm{\vec u+\vec v}^2$ grâce à l'identité
        $\norm{\vec u+\vec v}^2=\norm{\vec u}^2+\norm{\vec v}^2+2\,\vec u\cdot\vec v$.
\end{enumerate}

\exo{5} \textbf{Produit $(\vec u+\vec v)\cdot(\vec u-\vec v)$.}
On donne $\vec u\,(3;1)$ et $\vec v\,(1;2)$.
\begin{enumerate}[label=\textbf{\alph*)},leftmargin=2em,itemsep=2pt]
  \item Calculer $\norm{\vec u}^2-\norm{\vec v}^2$.
  \item En déduire $(\vec u+\vec v)\cdot(\vec u-\vec v)$, puis vérifier par un calcul direct.
\end{enumerate}

\end{document}
